\documentclass[10pt, letterpaper]{article}

% 宏包 (Packages):
\usepackage[
    ignoreheadfoot, % 设置页边距时不考虑页眉和页脚
    top=2 cm, % 正文内容与页面顶边缘的间距
    bottom=2 cm, % 正文内容与页面底边缘的间距
    left=2 cm, % 正文内容与页面左边缘的间距
    right=2 cm, % 正文内容与页面右边缘的间距
    footskip=1 cm, % 正文与页脚之间的间距
    % showframe % 用于调试显示页面框架 
]{geometry} % 用于调整页面几何结构
\usepackage[UTF8]{ctex} % 仅负责中文
\usepackage{titlesec} % 用于自定义章节标题
\usepackage{tabularx} % 用于制作带固定宽度列的表格
\usepackage{array} % tabularx 宏包的依赖项
\usepackage[dvipsnames]{xcolor} % 用于设置文本颜色
\definecolor{primaryColor}{RGB}{0, 0, 0} % 定义主色调
\usepackage{enumitem} % 用于自定义列表
\usepackage{fontawesome5} % 用于使用图标
\usepackage{amsmath} % 用于数学公式
\usepackage[
    pdftitle={dzn's CV},
    pdfauthor={dzn},
    pdfcreator={LaTeX with RenderCV},
    colorlinks=true,
    urlcolor=primaryColor
]{hyperref} % 用于链接、元数据和书签
\usepackage[pscoord]{eso-pic} % 用于在页面上放置浮动文本
\usepackage{calc} % 用于计算长度
\usepackage{bookmark} % 用于书签
\usepackage{lastpage} % 用于获取总页数
\usepackage{changepage} % 用于单列条目 (adjustwidth 环境)
\usepackage{paracol} % 用于双列和三列条目
\usepackage{ifthen} % 用于条件判断语句
\usepackage{needspace} % 用于避免在章节标题后紧接着直接分页
\usepackage{iftex} % 检查编译引擎是 pdflatex, xetex 还是 luatex

% 确保生成的 PDF 是机器可读的/可被 ATS 解析的：
\ifPDFTeX
    \input{glyphtounicode}
    \pdfgentounicode=1
    \usepackage[T1]{fontenc}
    \usepackage[utf8]{inputenc}
    \usepackage{lmodern}
\fi

\usepackage{charter}

% 一些设置 :
\raggedright
\AtBeginEnvironment{adjustwidth}{\partopsep0pt} % 移除 adjustwidth 环境前的间距
\pagestyle{empty} % 不显示页眉或页脚
\setcounter{secnumdepth}{0} % 取消章节编号
\setlength{\parindent}{0pt} % 取消缩进
\setlength{\topskip}{0pt} % 取消顶部间距
\setlength{\columnsep}{0.15cm} % 设置列间距
\pagenumbering{gobble} % 不显示页码

\titleformat{\section}{\needspace{4\baselineskip}\bfseries\large}{}{0pt}{}[\vspace{1pt}\titlerule]

\titlespacing{\section}{
    % 左侧间距:
    -1pt
}{
    % 顶部间距:
    0.3 cm
}{
    % 底部间距:
    0.2 cm
} % 章节标题间距设置

\renewcommand\labelitemi{$\vcenter{\hbox{\small$\bullet$}}$} % 自定义列表符号
\newenvironment{highlights}{
    \begin{itemize}[
        topsep=0.10 cm,
        parsep=0.10 cm,
        partopsep=0pt,
        itemsep=0pt,
        leftmargin=0 cm + 10pt
    ]
}{
    \end{itemize}
} % 高亮内容的新环境


\newenvironment{highlightsforbulletentries}{
    \begin{itemize}[
        topsep=0.10 cm,
        parsep=0.10 cm,
        partopsep=0pt,
        itemsep=0pt,
        leftmargin=10pt
    ]
}{
    \end{itemize}
} % 带项目符号条目的高亮内容新环境

\newenvironment{onecolentry}{
    \begin{adjustwidth}{
        0 cm + 0.00001 cm
    }{
        0 cm + 0.00001 cm
    }
}{
    \end{adjustwidth}
} % 单列条目的新环境

\newenvironment{twocolentry}[2][]{
    \onecolentry
    \def\secondColumn{#2}
    \setcolumnwidth{\fill, 4.5 cm}
    \begin{paracol}{2}
}{
    \switchcolumn \raggedleft \secondColumn
    \end{paracol}
    \endonecolentry
} % 双列条目的新环境

\newenvironment{threecolentry}[3][]{
    \onecolentry
    \def\thirdColumn{#3}
    \setcolumnwidth{, \fill, 4.5 cm}
    \begin{paracol}{3}
    {\raggedright #2} \switchcolumn
}{
    \switchcolumn \raggedleft \thirdColumn
    \end{paracol}
    \endonecolentry
} % 三列条目的新环境

\newenvironment{header}{
    \setlength{\topsep}{0pt}\par\kern\topsep\centering\linespread{1.5}
}{
    \par\kern\topsep
} % 页眉的新环境

\newcommand{\placelastupdatedtext}{% \placetextbox{<水平位置>}{<垂直位置>}{<内容>}
  \AddToShipoutPictureFG*{% 将内容添加到当前页面前景
    \put(
        \LenToUnit{\paperwidth-2 cm-0 cm+0.05cm},
        \LenToUnit{\paperheight-1.0 cm}
    ){\vtop{{\null}\makebox[0pt][c]{
        \small\color{gray}\textit{Last updated in September 2024}\hspace{\widthof{Last updated in September 2024}}
    }}}%
  }%
}%

% 将原始的 href 命令保存到一个新命令中：
\let\hrefWithoutArrow\href

% 用于外部链接的新命令：


\begin{document}
    \newcommand{\AND}{\unskip
        \cleaders\copy\ANDbox\hskip\wd\ANDbox
        \ignorespaces
    }
    \newsavebox\ANDbox
    \sbox\ANDbox{$|$}

    \begin{header}
        \fontsize{25 pt}{25 pt}\selectfont 戴正楠

        \vspace{5 pt}

        \normalsize
        \mbox{江苏南京}%
        \kern 5.0 pt%
        \AND%
        \kern 5.0 pt%
        \mbox{88888888888}%
        \kern 5.0 pt%
        \AND%
        \kern 5.0 pt%
        \mbox{\hrefWithoutArrow{https://github.com/zheshigewenti}{github.com/zheshigewenti}}%
    \end{header}

    \vspace{5 pt - 0.3 cm}




    



    \section{教育}



        
        \begin{twocolentry}{
            9月 2017 –  5月 2021
        }
            \textbf{金陵科技学院}\end{twocolentry}

        \vspace{0.10 cm}
        \begin{onecolentry}
            \begin{highlights}
                \item GPA: 3.9/5.0 
                \item \textbf{Coursework:} Computer Architecture, Comparison of Learning Algorithms, Computational Theory
            \end{highlights}
        \end{onecolentry}



    
    \section{经历}



        
        \begin{twocolentry}{
            7月 2021 – 6月 2022
        }
            \textbf{结构工程师}, AGG新能源\end{twocolentry}

        \vspace{0.10 cm}
        \begin{onecolentry}
            \begin{highlights}
                \item Reduced time to render user buddy lists by 75\% by implementing a prediction algorithm
                \item Integrated iChat with Spotlight Search by creating a tool to extract metadata from saved chat transcripts and provide metadata to a system-wide search database
                \item Redesigned chat file format and implemented backward compatibility for search
            \end{highlights}
        \end{onecolentry}


        \vspace{0.2 cm}

        \begin{twocolentry}{
            9月 2023 – 6月 2025
        }
            \textbf{监理工程师}, 中核工程咨询有限公司\end{twocolentry}

        \vspace{0.10 cm}
        \begin{onecolentry}
            \begin{highlights}
                \item Designed a UI for the VS open file switcher (Ctrl-Tab) and extended it to tool windows
                \item Created a service to provide gradient across VS and VS add-ins, optimizing its performance via caching
                \item Built an app to compute the similarity of all methods in a codebase, reducing the time from $\mathcal{O}(n^2)$ to $\mathcal{O}(n \log n)$
                \item Created a test case generation tool that creates random XML docs from XML Schema
                \item Automated the extraction and processing of large datasets from legacy systems using SQL and Perl scripts
            \end{highlights}
        \end{onecolentry}



    
 
    
    \section{项目}



        
        \begin{twocolentry}{
            \href{https://github.com/sinaatalay/rendercv}{github.com/name/repo}
        }
            \textbf{Multi-User Drawing Tool}\end{twocolentry}

        \vspace{0.10 cm}
        \begin{onecolentry}
            \begin{highlights}
                \item Developed an electronic classroom where multiple users can simultaneously view and draw on a "chalkboard" with each person's edits synchronized
                \item Tools Used: C++, MFC
            \end{highlights}
        \end{onecolentry}


        \vspace{0.2 cm}

        \begin{twocolentry}{
            \href{https://github.com/sinaatalay/rendercv}{github.com/name/repo}
        }
            \textbf{Synchronized Desktop Calendar}\end{twocolentry}

        \vspace{0.10 cm}
        \begin{onecolentry}
            \begin{highlights}
                \item Developed a desktop calendar with globally shared and synchronized calendars, allowing users to schedule meetings with other users
                \item Tools Used: C\#, .NET, SQL, XML
            \end{highlights}
        \end{onecolentry}


        \vspace{0.2 cm}

        \begin{twocolentry}{
            2002
        }
            \textbf{Custom Operating System}\end{twocolentry}

        \vspace{0.10 cm}
        \begin{onecolentry}
            \begin{highlights}
                \item Built a UNIX-style OS with a scheduler, file system, text editor, and calculator
                \item Tools Used: C
            \end{highlights}
        \end{onecolentry}



    
    \section{语言}



        
        \begin{onecolentry}
            \textbf{Languages:}  C, Nix, Lua, SQL, Python
        \end{onecolentry}

        \vspace{0.2 cm}

        \begin{onecolentry}
            \textbf{Technologies:} Neovim, Office,  Nixpkgs,  Interface Builder, CAD, Ltex
        \end{onecolentry}


    

\end{document}
